\chapter{The Circuit}
The next extension of the model is a very significant one, that introduces a new sector and two new assets.

The goal of this extension is to add a stylized financial sector inspired to Graziani circuit theory \parencite{passarella2024, caiani2017}: firms get loans from the financial sector to pay for wages and investments and repay them at the end of the period.

\begin{table}
    \centering
    \begin{tabular}{l|ccccc|l}
                 & $\mathcal{H}$             & $\mathcal{F}_\mathbf{C}$               & $\mathcal{F}_\mathbf{K}$               & $\mathcal{B}$    & $\mathcal{G}$    & $\Sigma$                   \\
        \midrule
        Money    & $+\mathbf{M}_\mathcal{H}$ & $+\mathbf{M}_{\mathcal{F}_\mathbf{C}}$ & $+\mathbf{M}_{\mathcal{F}_\mathbf{K}}$ & $-\mathbf{M}$    &                  & 0                          \\
        Loans    &                           & $-\mathbf{L}_{\mathcal{F}_\mathbf{C}}$ & $-\mathbf{L}_{\mathcal{F}_\mathbf{K}}$ & $+\mathbf{L}$    &                  & 0                          \\
        Bonds    &                           &                                        &                                        & $+\mathbf{B}$    & $-\mathbf{B}$    & 0                          \\
        \midrule
        Capital  &                           & $+p_\mathbf{K} \mathbf{K}_\mathbf{C}$  & $+p_\mathbf{K} \mathbf{K}_\mathbf{K}$  &                  &                  & $+p_\mathbf{K} \mathbf{K}$ \\
        \midrule
        $\Sigma$ & $+V_\mathcal{H}$          & $+V_{\mathcal{F}_\mathbf{C}}$          & $+V_{\mathcal{F}_\mathbf{K}}$          & $+V_\mathcal{B}$ & $+V_\mathcal{G}$ & $+p_\mathbf{K} \mathbf{K}$ \\
    \end{tabular}
    \caption{Balance sheet for the model with investments}
\end{table}

\begin{table}
    \centering
    \begin{tabular}{l|ccccc|l}

                               & $\mathcal{H}$                   & $\mathcal{F}_\mathbf{C}$                                                                                    & $\mathcal{F}_\mathbf{K}$                                                                                    & $\mathcal{B}$                                             & $\mathcal{G}$       & $\Sigma$ \\
        \midrule
        Consumption            & $-p C$                          & $+p C$                                                                                                      &                                                                                                             &                                                           &                     & 0        \\
        Gvt Spending           &                                 & $+p G$                                                                                                      &                                                                                                             &                                                           & $-p G$              & 0        \\
        Investments            &                                 & $-p_\mathbf{K} I_{\mathcal{F}_\mathbf{C}}$                                                                  & $+p_\mathbf{K} I_{\mathcal{F}_\mathbf{C}}$                                                                  &                                                           &                     & 0        \\
        Benefits               & $+UB$                           &                                                                                                             &                                                                                                             &                                                           & $-UB$               & 0        \\
        Wages                  & $+W$                            & $-W_{\mathcal{F}_\mathbf{C}}$                                                                               & $-W_{\mathcal{F}_\mathbf{K}}$                                                                               &                                                           &                     & 0        \\
        Profits                & $+P$                            & $-P_{\mathcal{F}_\mathbf{C}}$                                                                               & $-P_{\mathcal{F}_\mathbf{K}}$                                                                               & $-P_\mathcal{B}$                                          &                     & 0        \\
        Taxes                  &
        $-T$                   &                                 &                                                                                                             &                                                                                                             & $+T$                                                      & 0                              \\
        \midrule
        $\mathbf{L}$ interests &                                 & $-r_\mathbf{L} (\Delta^+\mathbf{L}_{\mathcal{F}_\mathbf{C}} + {\mathbf{L}_{\mathcal{F}_\mathbf{C}}}_{t-1})$ & $-r_\mathbf{L} (\Delta^+\mathbf{L}_{\mathcal{F}_\mathbf{K}} + {\mathbf{L}_{\mathcal{F}_\mathbf{K}}}_{t-1})$ & $+r_\mathbf{L} (\Delta^+\mathbf{L} + {\mathbf{L}}_{t-1})$ &                     & 0        \\
        $\mathbf{B}$ interests &
                               &                                 &                                                                                                             & $+r_\mathbf{B}\mathbf{B}$                                                                                   & $-r_\mathbf{B}\mathbf{B}$                                 & 0                              \\
        \midrule
        $\Delta$ Money         & $-\Delta\mathbf{M}_\mathcal{H}$ & $-\Delta\mathbf{M}_{\mathcal{F}_\mathbf{C}}$                                                                & $-\Delta\mathbf{M}_{\mathcal{F}_\mathbf{K}}$                                                                & $+\Delta\mathbf{M}$                                       &                     & 0        \\
        $\Delta$ Loans         &                                 & $+\Delta\mathbf{L}_{\mathcal{F}_\mathbf{C}}$                                                                & $+\Delta\mathbf{L}_{\mathcal{F}_\mathbf{K}}$                                                                & $-\Delta\mathbf{L}_{\mathcal{B}}$                         &                     & 0        \\
        $\Delta$ Bonds         &                                 &                                                                                                             &                                                                                                             & $-\Delta\mathbf{B}$                                       & $+\Delta\mathbf{B}$ & 0        \\
        \midrule
        $\Sigma$               & 0                               & 0                                                                                                           & 0                                                                                                           & 0                                                         & 0                   & 0        \\
    \end{tabular}
    \caption{Transactions and Flows matrix for the model with investments}
\end{table}

The matrices of the models change as below: money is now issued by the bank as deposits and backed by government bonds (which can serve as reserves if the sign switches), while firms have access to credit through the loans issued by the banking sector.
Both loans and bonds are interests bearing.

\section{The aggregate version}
\subsection{The behavioural equations}
Firms receive loans from the bank to cover for wages and investment $$\Delta^+ \mathbf{L}_{\mathcal{F}_\mathbf{C}} = W_{\mathcal{F}_\mathbf{C}} + p_\mathbf{K} I_{\mathcal{F}_\mathbf{C}}$$ $$\Delta^+\mathbf{L}_{\mathcal{F}_\mathbf{K}} = W_{\mathcal{F}_\mathbf{K}}$$
At the end of the period the stock of money of the previous period and the income of the current one are used to repay the loans $$\mathbf{L}_{\mathcal{F}_\mathbf{C}} = {\mathbf{L}_{\mathcal{F}_\mathbf{C}}}_{t-1} + \Delta^+\mathbf{L}_{\mathcal{F}_\mathbf{C}} - {\mathbf{M}_{\mathcal{F}_\mathbf{C}}} - p_\mathbf{C} (C + G)$$ $$\mathbf{L}_{\mathcal{F}_\mathbf{K}} = {\mathbf{L}_{\mathcal{F}_\mathbf{K}}}_{t-1} + \Delta^+\mathbf{L}_{\mathcal{F}_\mathbf{K}} - {\mathbf{M}_{\mathcal{F}_\mathbf{K}}} - p_\mathbf{K} I_{\mathcal{F}_\mathbf{C}}$$

Then interest are paid on maximum amount of loans inside the period $\mathbf{L}_{t-1} + \Delta^+\mathbf{L}$

Bonds are issued by the government to finance the deficit and interests are paid on them; the accounting at the end of the period is $$\mathbf{B} = \frac{\mathbf{B}_{t-1} - (T - p_\mathbf{C}G - UB)}{1-r_\mathbf{B}}$$
where $r$ denotes the interest rate.

The interest rate on bond is determined with a Taylor-like rule depending on inflation ($i$), unemployment ($u$) and the capacity utilization as $$r_\mathbf{B} = i_{t-1} + \alpha_i (i_{i-1} - i^T) + \alpha_{cu} ({cu}_{t-1} - {cu}^T) - \alpha_u (u_{t-1} - u^T)$$ and the interest rate on loan is determined as a fixed premium on the interest rate on bonds.

The bank aims to keep a target wealth to loan ratio, which leads to the equation for the distributed profits (see appendix \ref{app:eq}) $$ P_\mathcal{B}  = \frac{1-r_\mathbf{B}}{1-r_\mathbf{B}+r_\mathbf{B}\tau_P}(\mathbf{V}_{t-1} - {cr}^T \mathbf{L} + r_\mathbf{L} (\Delta^+\mathbf{L} + \mathbf{L}_{t-1}) + \frac{r_\mathbf{B}}{1-r_\mathbf{B}}(\mathbf{B}_{t-1} + p_\mathbf{C} G + UB - \tau_C C - \tau_W W - \tau_P (P_{\mathcal{F}_\mathbf{C}} + P_{\mathcal{F}_\mathbf{K}})))$$

\subsection{Results}
\begin{figure}
    \centering
    \includegraphics[width=\textwidth]{04_sankey}
    \caption{Sankey diagram of flow at the steady state}
    \label{fig:04_san}
\end{figure}

\begin{figure}
    \centering
    \begin{subfigure}[t]{0.49\textwidth}
        \centering
        \includegraphics[width=\textwidth]{04_L}
        \caption{Stock of Loans}
        \label{fig:04_L}
    \end{subfigure}%
    ~
    \begin{subfigure}[t]{0.49\textwidth}
        \centering
        \includegraphics[width=\textwidth]{04_P}
        \caption{Distributed profits}
        \label{fig:04_P}
    \end{subfigure}
    \caption{Dynamics of the model with the circuit}
\end{figure}

The sankey diagram (figure \ref{fig:04_san}) looks significantly different, but the dynamics of the other variables is the same of the previous version.

The profits are mostly concentrated in the consumption goods sector (figure \ref{fig:04_P}) and the circuit closes entirely, not leaving a stock of loans at the end of the period (figure \ref{fig:04_L}).

\subsection{Exogenous government spending}
The model with exogenous government spending reach full-employment and a stable dynamics when the financial sector is added.
This is clearly an effect of the endogeineity of money, which is now created without limits by the financial sector, being able to accommodate financially any growth of demand.

To keep the demand constrained and the model to not reach full employment, it is necessary to also constrain the credit given to the firms.
We impose as the restriction on the credit lines for each sector that $$\Delta^+\mathbf{L} \le V_{t-1}$$ representing the necessity of a collateral good required by the bank.

Already in the previous versions the markup rate of the capital sector decline until it is not sufficient to cover for the cost of self-investment.
Moreover, his version adds the interest on the loans taken to start the circuit as a cost of production, which is not taken in account in determining the prices.
It is possible to either adjust the price equation or to fix a minimum value for the mark-up. The second option has been chosen.

By adding these further restrictions the model preserve its features: it still does not reach full employment, the demand for consumption goods is fully satisfied, and it presents cyclical behaviour.

\section{The matrix-based agent-based extension}
By adding another sector, it is necessary to choose with how many agents it should be represented.
For realism, financial sector can be populated by competing banks, financial firms that manage savings and trade in bonds and equities.
But this would increase the complexity of the model (by adding many new interactions).
As discussed in the chapter \ref{ch:complexity} we need to identify which level of detail better suits the needs.
The model is rooted in the production side of the economy, labour and goods are the core elements we are investigating.
To tame the complexity of the model we choose to consider the financial sector as a very stylized aggregate sector rather than a detailed sector with multiple agents.

The agent based follows closely the aggregate version with endogenous government spending, including the absence of constrains on credit.
The sequence of the events of the model is updated as:
\begin{enumerate}[label=\textbf{\Alph*}]
    \item Households determinate the desired consumption (expressed as units of good) based on the wealth and the disposable income of the previous period
    \item Government determinate the desired consumption (expressed as units of good) based on the target deficit computed on the accounting of the previous period
    \item Investments are ordered by the consumption firms to the capital ones
    \item Firms fire and hire workers in order to match the labour input required to match the desired consumption
    \item The realized output is determined for both capital and consumption goods
    \item Wages and unemployment benefits are paid
    \item Prices are set
    \item Consumption goods are sold to households and the government
    \item Capital goods are sold to consumption firms
    \item Loans are written down to cover wages and investments
    \item Loans are repaid in their principal and interests
    \item Profits are distributed
    \item Taxes are paid
    \item Accounting variables are updated, including government bonds' interest
\end{enumerate}

The only relevant difference is a consequence of the heterogeneity of the firms: as a consequence of the random matching, some firms accumulate enough debt to reach negative net wealth.
Those firms are considered bankrupted, their stock of loans is cancelled (the loss is absorbed by the bank) and they restart the production the following period with their stock of capital.

\subsection{Results}
\begin{figure}
    \centering
    \includegraphics[width=\textwidth]{04AB_L}
    \caption{Dynamics of the matrix-based AB model with the circuit}
    \label{fig:04AB_L}
\end{figure}

In figure \ref{fig:04AB_L} we can observe the spike in the stock of loans held by capital firms before the bankruptcies. Additionally, profits are negligible in the capital firms sector and are distributed mostly by the banking one and the consumption firms one.

Finally, the fat-tail distribution of wealth among households becomes even more skewed, resembling a one-gets-all dynamics.

\section{The OOP agent-based extension}
As in previous iterations, the OOP version is less regular in its behaviour, here creating short term fluctations in debt dynamics which however smooth out the medium term dynamics (figure \ref{fig:04ABO_L}).

\begin{figure}
    \centering
    \includegraphics[width=\textwidth]{04ABO_L}
    \caption{Dynamics of the OOP AB model with the circuit}
    \label{fig:04ABO_L}
\end{figure}